\documentclass[11pt]{article}

\usepackage[margin=1in]{geometry}
\usepackage{amsmath,amssymb,amsthm}
\usepackage{booktabs}
\usepackage{hyperref}
\usepackage{microtype}

\newtheorem{proposition}{Proposition}
\newtheorem{corollary}[proposition]{Corollary}
\theoremstyle{definition}
\newtheorem{definition}[proposition]{Definition}

\title{Structural Exclusion of Shift-Based Cohorts in Panel Scheduling}
\author{Mihir Bommisetty}
\date{}

\begin{document}
\maketitle

\begin{abstract}
Panel assessment processes---hiring loops, promotion boards, certification reviews---are
typically scheduled as a block of consecutive sessions followed by a synchronous debrief.
We show that this format structurally excludes cohorts whose availability is confined to
short windows adjacent to a work shift, independently of their number, qualification, or
willingness to participate. We formalise \emph{cohort reachability} for a population with
heterogeneous availability, and prove that exclusion arises from two distinct and
independent mechanisms: \emph{contiguity}, which requires a run of consecutively staffed
slots at least as long as the panel, and \emph{simultaneity}, which requires the entire
panel to share a common free slot for debrief. We show that the second implies that
cohorts with pairwise-disjoint availability can never share a panel under any synchronous
debrief policy, and hence that such cohorts are served only by single-cohort panels. A
constraint-programming implementation at a scale of $10^5$ assessors confirms the
analysis: of three candidate interventions, the two that address contiguity or debrief
timing alone yield no improvement, while relaxing the contiguity requirement---the only
one that admits single-cohort panels---moves the affected cohort from $0\%$ to
approximately $15\%$ of panel seats at no cost to the proportion of panels filled.
Results are obtained on synthetic populations; the mechanisms are exact, the magnitudes
are not validated against field data.
\end{abstract}

\section{Introduction}

A substantial fraction of the workforce in manufacturing, logistics, healthcare, and
retail does not work at a desk. These employees are typically assigned to rotating shifts,
lack individual corporate accounts, and have narrow discretionary time. Practitioner
literature on ``deskless'' or ``frontline'' workers observes that processes designed
around calendar software and email systematically under-reach them.

This paper examines one such process: panel assessment, in which several assessors each
conduct a session with a subject, followed by a joint debrief. The dominant format
schedules all sessions consecutively on a single day. We ask whether this format is
merely inconvenient for shift workers, or whether it excludes them as a matter of
combinatorial structure.

We find the latter, and identify two separate mechanisms. Our contributions are:

\begin{enumerate}
  \item A formalisation of \emph{cohort reachability} under a scheduling protocol
        (Section~\ref{sec:model}).
  \item A characterisation of reachability under contiguous protocols in terms of the
        connectivity of the staffed set, including the observation that a single
        unstaffed slot can exclude a cohort no member of which is available at that slot
        (Propositions~\ref{prop:cont}--\ref{prop:sever}).
  \item A stronger exclusion result for synchronous debrief: panels are schedulable only
        if all assessors share a free slot, so cohorts with disjoint availability can
        never co-serve, and are reachable only via single-cohort panels
        (Propositions~\ref{prop:debrief}--\ref{prop:disjoint}).
  \item A computational study at $10^5$ assessors evaluating three interventions, two of
        which the theory predicts---correctly---will fail (Section~\ref{sec:exp}).
\end{enumerate}

\section{Related work}

Constraint programming for scheduling is a mature field, with established formulations for
job-shop scheduling, nurse rostering, timetabling, and vehicle routing with shift
constraints. Our formulation uses standard techniques and we claim no methodological
novelty there.

Work on scheduling equity has largely addressed the fairness of shift assignment
itself---who works nights, how rest is distributed, how preferences are honoured. To our
knowledge the complementary question, whether a shift structure excludes workers from
\emph{other} organisational processes, has not been treated formally. The practitioner
literature on frontline-worker exclusion is qualitative and does not model the mechanism.

\section{Model}
\label{sec:model}

Let $T=\{1,\dots,n\}$ be a discretised horizon of slots. Each assessor $p$ in a population
$P$ has an availability set $A_p \subseteq T$. A \emph{cohort} $C$ is a maximal set of
assessors sharing an availability pattern, written $A_C$; in our setting cohorts
correspond to shifts.

A process requires $k$ sessions. A \emph{protocol} $\Pi \subseteq T^k$ is the set of
permitted session-time tuples. We consider

\[
\Pi_{\mathrm{cont}} = \{(t,t+1,\dots,t+k-1)\},
\qquad
\Pi_{\mathrm{split}} = \{(t_1,\dots,t_k) : t_i \text{ distinct}\}.
\]

Let $N(t)=\{p : t \in A_p\}$ and let $S=\{t : N(t)\neq\emptyset\}$ denote the
\emph{staffed set}. A schedule $s\in\Pi$ is \emph{viable} if $N(s_i)\neq\emptyset$ for
all $i$.

\begin{definition}[Reachability]
Assessor $p$ is \emph{reachable} under $\Pi$ if there is a viable $s\in\Pi$ and an index
$i$ with $s_i \in A_p$. A cohort is \emph{excluded} if no member is reachable.
\end{definition}

Reachability is necessary but not sufficient for participation: it states that an
assignment is possible, not that it occurs, and says nothing about share.

\section{Contiguity}

Write $S$ as a disjoint union of maximal runs $R_1,\dots,R_m$ of consecutive slots.

\begin{proposition}
\label{prop:cont}
Under $\Pi_{\mathrm{cont}}$, assessor $p$ is reachable if and only if there is a run $R_j$
with $|R_j| \ge k$ and $A_p \cap R_j \neq \emptyset$.
\end{proposition}

\begin{proof}
($\Leftarrow$) Given $t^\ast \in A_p \cap R_j$ with $|R_j|\ge k$, select $k$ consecutive
slots within $R_j$ containing $t^\ast$. Every such slot lies in $S$, so the schedule is
viable.
($\Rightarrow$) A viable contiguous schedule occupies $k$ consecutive slots of $S$, which
lie in a common maximal run $R_j$; hence $|R_j|\ge k$, and if $p$ fills a session then
$A_p$ meets $R_j$.
\end{proof}

The content lies in the quantifier: reachability depends on the run containing the
cohort's window, not on the length of that window.

\begin{corollary}
If $A_C \subseteq R_j$ with $|R_j| < k$, cohort $C$ is excluded.
\end{corollary}

\begin{proposition}[Severance]
\label{prop:sever}
There exist populations and slots $t^\dagger$ such that some cohort $C$ is reachable with
respect to $S$ and excluded with respect to $S\setminus\{t^\dagger\}$, despite
$t^\dagger \notin A_C$.
\end{proposition}

\begin{proof}
Take $k=4$, $A_C=\{18,19\}$ and $A_D=\{14,15,16,17\}$. Then $S=\{14,\dots,19\}$ is a run
of length $6$, and by Proposition~\ref{prop:cont} members of $C$ are reachable via the
window $\{16,17,18,19\}$. Removing slot $17$ splits $S$ into runs of length $3$ and $2$,
both shorter than $k$; no viable contiguous schedule exists and $C$ is excluded.
\end{proof}

Proposition~\ref{prop:sever} is the reason we distinguish reachability from window
length. In our setting slot $17$ corresponds to shift changeover, when one crew has left
and the next has not started.

\section{Simultaneity}

Viability as defined asks only that each session slot be staffed by \emph{someone}. Panel
processes typically add a debrief in which all $k$ assessors participate together.

\begin{proposition}
\label{prop:debrief}
Let $d(s)$ denote the debrief slot for schedule $s$. A panel $(p_1,\dots,p_k)$ is
debrief-feasible only if $d(s) \in \bigcap_i A_{p_i}$.
\end{proposition}

\begin{proposition}
\label{prop:disjoint}
Under any synchronous debrief policy, a panel is schedulable only if
$\bigcap_i A_{p_i} \neq \emptyset$. Consequently cohorts with pairwise-disjoint
availability can never share a panel, at any debrief time.
\end{proposition}

\begin{corollary}
\label{cor:single}
A cohort whose availability is disjoint from that of every other cohort is served only by
\emph{single-cohort} panels, in which all $k$ assessors are drawn from that cohort.
\end{corollary}

Corollary~\ref{cor:single} is the operative result. It explains why interventions that
adjust debrief \emph{timing} cannot help: the obstruction is disjointness, not schedule
placement.

\section{Split protocols}

\begin{proposition}
\label{prop:split}
Under $\Pi_{\mathrm{split}}$, assessor $p$ is reachable if and only if
$A_p \neq \emptyset$ and $|S| \ge k$.
\end{proposition}

Combining Proposition~\ref{prop:split} with Corollary~\ref{cor:single}: a single-cohort
panel drawn from cohort $C$ is schedulable under $\Pi_{\mathrm{split}}$ whenever
$|A_C| \ge k$ across the horizon, even if $A_C$ contains no run of length $k$. Under
$\Pi_{\mathrm{cont}}$ it requires such a run. This is the precise sense in which
relaxing contiguity, and only that, admits the excluded cohort.

\section{Computational study}
\label{sec:exp}

\subsection{Setup}

We generate synthetic populations of $10^5$ assessors distributed over six sites and six
functions. Corporate assessors are available 09:00--17:00 on weekdays; day-shift crews
18:00--20:00 (post-shift); swing and night crews in short pre-shift windows; weekend crews
Friday to Sunday. Assessor levels are drawn from a fixed distribution, and a per-assessor
weekly cap limits load.

Panels are assigned by a constraint program, sharded by site, solved with CP-SAT. The
model enforces exactly one assessor per session, no assessor in two places at once, no
assessor twice on a panel, a synchronous debrief with all panelists free, at least one
senior assessor, at least one assessor outside the subject's function, and per-assessor
caps. The objective maximises panels filled with a secondary term minimising the maximum
per-assessor load. Two implementation details are load-bearing at this scale: a
precomputed availability index keyed by (site, day, slot), and a capped candidate
shortlist per session---without the latter the model exceeds available memory.

\subsection{Interventions}

We evaluate three interventions against the standard contiguous protocol
(Table~\ref{tab:interventions}). The first staffs the changeover slot, addressing
Proposition~\ref{prop:sever}. The second relaxes debrief to any slot within 24 hours of
the final session. The third replaces $\Pi_{\mathrm{cont}}$ with $\Pi_{\mathrm{split}}$.

\begin{table}[h]
\centering
\begin{tabular}{lcc}
\toprule
Intervention & Day-shift seats & Panels filled \\
\midrule
None (baseline)                       & $0\%$      & $74.6\%$ \\
Staff changeover slot                 & $0\%$      & $60.9\%$ \\
Relax debrief to any slot within 24h  & $0\%$      & $60.9\%$ \\
Split sessions across days            & $\sim15\%$ & $73.7\%$ \\
\bottomrule
\end{tabular}
\caption{Day-shift share of panel seats under each intervention. Figures vary by
approximately one percentage point across runs owing to solver non-determinism.}
\label{tab:interventions}
\end{table}

The first two interventions fail, as Corollary~\ref{cor:single} predicts: neither admits a
single-cohort day-shift panel. Staffing the changeover slot does restore reachability in
the contiguity sense---with the debrief constraint removed, day-shift share rises to
$4.8\%$---but the simultaneity constraint re-excludes the cohort.

The third succeeds, and inspection of the resulting assignments confirms the predicted
structure: day-shift assessors appear on panels composed entirely of day-shift assessors,
with sessions distributed across days within their 18:00--20:00 windows.

\subsection{Cost}

Splitting sessions does not reduce the proportion of panels filled ($74.6\%$ versus
$73.7\%$, within run-to-run variation) and does not increase maximum per-assessor load.
We do not claim this is general; it is a property of the populations tested.

\section{Limitations}

All results are obtained on synthetic populations. Shift windows, cohort proportions, and
assessor level distributions are modelled rather than observed. The structural results of
Sections 4--6 are exact and independent of these choices; the magnitudes in
Table~\ref{tab:interventions} are not, and we make no claim that $15\%$ would be recovered
in any particular organisation.

We model reachability and observed share, but do not characterise the relationship between
them. Proposition~\ref{prop:split} establishes that the day shift becomes reachable; it
does not predict $15\%$. Closing that gap is the natural next question.

Finally, we treat availability as exogenous and binary. Real assessors have preferences,
fatigue, and partial availability, none of which we model.

\section{Conclusion}

Panel processes exclude shift-based cohorts through two independent mechanisms:
contiguity, which requires a sufficiently long run of staffed slots, and simultaneity,
which requires the panel to share a free slot. Interventions addressing either mechanism
alone are ineffective when both bind. Relaxing contiguity succeeds because it is the only
intervention among those tested that permits single-cohort panels, which
Corollary~\ref{cor:single} shows to be the sole available structure for a cohort whose
availability is disjoint from all others.

The framing generalises beyond hiring to any organisational process requiring several
participants over a contiguous block: safety committees, certification boards, promotion
panels. Where participation is treated as a matter of willingness or qualification, it may
instead be a matter of schedule structure.

\end{document}
